\documentclass[a4paper]{arrowhead}

\usepackage[yyyymmdd]{datetime}
\usepackage{etoolbox}
\usepackage[utf8]{inputenc}
\usepackage[T1]{fontenc}
\usepackage{multirow}
\usepackage{longtable}
\usepackage{tikz}
\usetikzlibrary{positioning,shapes.multipart,arrows.meta}

\renewcommand{\dateseparator}{-}
\setlength{\emergencystretch}{3em}
\lstset{language=HTTP,breakatwhitespace=false,columns=fullflexible}

%% Special references
\newcommand{\mref}[2][]{{\textcolor{ArrowheadPurple}{\hyperref[sec:model:#2]{\ifstrempty{#1}{#2}{#1}}}}}
\newcommand{\oref}[1]{{\textcolor{ArrowheadBlue}{\hyperref[sec:operations:#1]{#1}}}}
\newcommand{\pdef}[1]{{\textcolor{ArrowheadGrey}{#1\label{sec:model:primitives:#1}\label{sec:model:primitives:#1s}\label{sec:model:primitives:#1es}}}}
\newcommand{\pref}[1]{{\textcolor{ArrowheadGrey}{\hyperref[sec:model:primitives:#1]{#1}}}}

% as template_sd.tex; the input and output arguments are passed already formatted (\mref or \pref)
\newrobustcmd\osubsection[3]{
  \addtocounter{subsection}{1}
  \addcontentsline{toc}{subsection}{\protect\numberline{\thesubsection}operation \textcolor{ArrowheadBlue}{#1}}
  \subsection*{
    \thesubsection\quad
    operation
    \textcolor{ArrowheadBlue}{#1}
    (\notblank{#2}{#2}{})
    \notblank{#3}{: #3}{}
  }
  \label{sec:operations:#1}
}
\newrobustcmd\msubsection[2]{
  \addtocounter{subsection}{1}
  \addcontentsline{toc}{subsection}{\protect\numberline{\thesubsection}#1 \textcolor{ArrowheadPurple}{#2}}
  \subsection*{\thesubsection\quad#1 \textcolor{ArrowheadPurple}{#2}}
  \label{sec:model:#2} \label{sec:model:#2s} \label{sec:model:#2es}
}
%%

\begin{document}
\sloppy

%% Arrowhead Document Properties
\ArrowheadTitle{onboardingManagement}
\ArrowheadServiceID{onboardingManagement}
\ArrowheadType{Service Description}
\ArrowheadTypeShort{SD}
\ArrowheadVersion{5.2.2}
\ArrowheadDate{2026-10-30}
\ArrowheadAuthor{Jerker Delsing}
\ArrowheadStatus{RELEASE CANDIDATE}
\ArrowheadContact{jerker.delsing@ltu.se}
\ArrowheadFooter{\href{www.arrowhead.eu}{www.arrowhead.eu}}
\ArrowheadSetup
%%

%% Front Page
\begin{center}
  \vspace*{1cm}
  \huge{\arrowtitle}

  \vspace*{0.2cm}
  \LARGE{\arrowtype}
  \vspace*{1cm}

  \Large{Service ID: \textit{"\arrowid"}}

  \vspace*{\fill}

  % Front Page Image

  \vspace*{1cm}
  \vspace*{\fill}

  % Front Page Abstract
  \begin{abstract}
    This document describes the \texttt{onboarding\allowbreak{}Management} service, version 1.0.0, produced by the Onboarding support system of Arrowhead 5.2.2.
  \end{abstract}

  \vspace*{1cm}
\end{center}

\newpage
%%

%% Table of Contents
\tableofcontents
\newpage
%%

\section{Overview}
\label{sec:overview}

This document describes the \texttt{onboarding\allowbreak{}Management} service, which provides operator approval and management of device enrollments.
It is produced by the Onboarding system \cite{onb-sysd} in service version 1.0.0.

The rest of this document is organized as follows.
In Section \ref{sec:operations}, we describe the abstract message operations provided by the service.
In Section \ref{sec:model}, we end the document by presenting the data types used by the mentioned operations.
The concrete interface is specified in the HTTP/TLS/JSON IDD \cite{onb-onboardingManagement-http-idd}.

\subsection{Significant Prior Art}

Replaces the v4 shared-secret enrollment.

\subsection{How This Service Is Meant to Be Used}

An operator approves a device with \texttt{approve-\allowbreak{}device} (bootstrap-certificate fingerprint and the systems it may onboard), inspects enrollments, off-boards a device or removes the enrollment.

\subsection{Important Delimitations and Dependencies}
\label{sec:delimitations}

Needs an operator identity token bound to the operator certificate. An approval lists 1 to 32 system names for the device; reserved names (core, support and operator systems, configurable with \texttt{ONBOARDING\_\allowbreak{}RESERVED\_\allowbreak{}NAMES}) are refused.

\newpage

\section{Service Interface}
\label{sec:operations}

This section describes the interfaces to the \texttt{onboarding\allowbreak{}Management} service.
In particular, each subsection names an abstract operation, an input type and an output type, in that order.
The input type is named inside parentheses, while the output type is preceded by a colon.
Input and output types are only denoted when accepted or returned, respectively, by the interface in question.

All abstract data types named in this section are defined in Section \ref{sec:model}.

\begin{figure}[ht!]
  \centering
  \begin{tikzpicture}[font=\small]
    \node[draw=ArrowheadBlue, very thick, rectangle split, rectangle split parts=2, align=center,
          rectangle split part fill={ArrowheadBlue!12,white}, text width=9cm, inner sep=6pt] (b)
      {\guillemotleft{}Service\guillemotright{}\\\textbf{onboardingManagement}
        \nodepart{second} \texttt{query-\allowbreak{}enrollments}\\ \texttt{get-\allowbreak{}enrollment}\\ \texttt{approve-\allowbreak{}device}\\ \texttt{offboard-\allowbreak{}device}\\ \texttt{remove-\allowbreak{}enrollment}};
  \end{tikzpicture}
  \caption{SysML-style block of the \texttt{onboarding\allowbreak{}Management} service and its operations.}
  \label{fig:SD}
\end{figure}

The following interface operations are available.

\osubsection{query-enrollments}{}{\mref{EnrollmentList}}

Returns all enrollments.

\osubsection{get-enrollment}{\pref{Name}}{\mref{EnrollmentView}}

Returns the enrollment of a device.

\osubsection{approve-device}{\mref{EnrollmentRequest}}{\mref{EnrollmentView}}

Approves a device: bootstrap-certificate fingerprint and the systems it may onboard.

\osubsection{offboard-device}{\pref{Name}}{\mref{OffboardingReport}}

Off-boards a device on behalf of the operator (as \texttt{onboarding.\allowbreak{}offboard-\allowbreak{}device}).

\osubsection{remove-enrollment}{\pref{Name}}{\mref{OffboardingReport}}

Off-boards the device and removes its enrollment.

\newpage

\section{Information Model}
\label{sec:model}

Here, all data objects that can be part of the \texttt{onboarding\allowbreak{}Management} service are listed in alphabetic order.
Note that each subsection, which describes one type of object, begins with the \textit{struct} keyword, which is used to denote a collection of named fields, each with its own data type.
As a complement to the explicitly defined types in this section, there is also a list of implicit primitive types in Section \ref{sec:model:primitives}, which are used to represent things like hashes and identifiers.

\msubsection{struct}{EnrollmentList}

List of enrollments.


\noindent\begin{tabularx}{\textwidth}{| >{\raggedright\arraybackslash}p{3.9cm} | >{\raggedright\arraybackslash\small}p{5.3cm} | X |} \hline
\rowcolor{gray!33} Field & Type & Description \\ \hline
entries & \pref{List}$<$\allowbreak{}\mref{EnrollmentView}$>$ & Matching records of the current page. \\ \hline
count & \pref{Number} & Number of enrollments. \\ \hline
\end{tabularx}



\msubsection{struct}{EnrollmentRequest}

Operator approval of a device.


\noindent\begin{tabularx}{\textwidth}{| >{\raggedright\arraybackslash}p{3.9cm} | >{\raggedright\arraybackslash\small}p{5.3cm} | X |} \hline
\rowcolor{gray!33} Field & Type & Description \\ \hline
bootstrap\allowbreak{}Fingerprint & \pref{String} & SHA-256 fingerprint of the approved bootstrap certificate. \\ \hline
systems & \pref{List}$<$\allowbreak{}\pref{String}$>$ & Systems the device may onboard. \\ \hline
\end{tabularx}



\msubsection{struct}{EnrollmentView}

An enrollment (approved device).


\noindent\begin{tabularx}{\textwidth}{| >{\raggedright\arraybackslash}p{3.9cm} | >{\raggedright\arraybackslash\small}p{5.3cm} | X |} \hline
\rowcolor{gray!33} Field & Type & Description \\ \hline
device\allowbreak{}Name & \pref{Name} & Name of the device hosting the system. \\ \hline
bootstrap\allowbreak{}Fingerprint & \pref{String} & SHA-256 fingerprint of the approved bootstrap certificate. \\ \hline
approved\allowbreak{}Systems & \pref{List}$<$\allowbreak{}\pref{String}$>$ & Systems the operator approved for the device. \\ \hline
state & \pref{String} & Enrollment state: \texttt{APPROVED}, \texttt{DEVICE\_\allowbreak{}REGISTERED} or \texttt{OFFBOARDED}. \\ \hline
device\allowbreak{}Certificate\allowbreak{}Id & \pref{Number} & Certificate\allowbreak{}Authority record of the device certificate. \\ \hline
device\allowbreak{}Certificate\allowbreak{}Expires\allowbreak{}At & \pref{DateTime} & Expiry of the device certificate. \\ \hline
systems & \pref{List}$<$\allowbreak{}\mref{SystemEnrollmentView}$>$ & Systems onboarded for the device. \\ \hline
approved\allowbreak{}By & \pref{String} & Operator who approved the device. \\ \hline
approved\allowbreak{}At & \pref{DateTime} & Time of the approval. \\ \hline
updated\allowbreak{}At & \pref{DateTime} & UTC time of the last change. \\ \hline
last\allowbreak{}Message & \pref{String} & Last status message of the enrollment. \\ \hline
\end{tabularx}



\msubsection{struct}{OffboardingReport}

Outcome of an off-boarding.


\noindent\begin{tabularx}{\textwidth}{| >{\raggedright\arraybackslash}p{3.9cm} | >{\raggedright\arraybackslash\small}p{5.3cm} | X |} \hline
\rowcolor{gray!33} Field & Type & Description \\ \hline
device\allowbreak{}Name & \pref{Name} & The off-boarded device. \\ \hline
completed & \pref{List}$<$\allowbreak{}\pref{String}$>$ & Off-boarding steps that succeeded. \\ \hline
failed & \pref{List}$<$\allowbreak{}\pref{String}$>$ & Off-boarding steps that failed. \\ \hline
\end{tabularx}



\msubsection{struct}{SystemEnrollmentView}

A system onboarded for a device.


\noindent\begin{tabularx}{\textwidth}{| >{\raggedright\arraybackslash}p{3.9cm} | >{\raggedright\arraybackslash\small}p{5.3cm} | X |} \hline
\rowcolor{gray!33} Field & Type & Description \\ \hline
system\allowbreak{}Name & \pref{Name} & Name of the system, PascalCase (e.g. \texttt{TemperatureSensor}). \\ \hline
certificate\allowbreak{}Id & \pref{Number} & Certificate\allowbreak{}Authority record of the system certificate. \\ \hline
certificate\allowbreak{}Expires\allowbreak{}At & \pref{DateTime} & Expiry of the system certificate. \\ \hline
issued\allowbreak{}At & \pref{DateTime} & Time the certificate was issued. \\ \hline
\end{tabularx}



\subsection{Primitives}
\label{sec:model:primitives}

Types and structures mentioned throughout this document that are assumed to be available to implementations of this service.
The concrete interpretations of each of these types and structures must be provided by any IDD document claiming to implement this service.

\noindent\begin{tabularx}{\textwidth}{| p{4.6cm} | X |} \hline
\rowcolor{gray!33} Type & Description \\ \hline
\pdef{DateTime} & Pinpoints a specific moment in time (UTC). \\ \hline
\pdef{List} & An \textit{array} of a known number of items, each having type A. \\ \hline
\pdef{Name} & A string identifier that is intended to be both human and machine-readable. \\ \hline
\pdef{Number} & An integer number. \\ \hline
\pdef{String} & A chain of characters. \\ \hline
\end{tabularx}

\newpage

\bibliographystyle{IEEEtran}
\bibliography{bibliography}

\newpage

\section{Revision History}
\subsection{Amendments}

\noindent\begin{tabularx}{\textwidth}{| p{1cm} | p{3cm} | p{2cm} | X | p{4cm} |} \hline
\rowcolor{gray!33} No. & Date & Version & Subject of Amendments & Author \\ \hline

1 & 2026-10-22 & \arrowversion & Document created from the Arrowhead LaTeX templates for the 5.2.2 Java Spring implementation, including the local extensions of this workspace. & Jerker Delsing \\ \hline
2 & 2026-10-30 & \arrowversion & Checked against the primary Java source tree (Eclipse-Arrowhead-java): approval limits (1 to 32 systems, reserved names) added. & Claude Opus 5.5 (draft) \\ \hline
\end{tabularx}

\subsection{Quality Assurance}

\noindent\begin{tabularx}{\textwidth}{| p{1cm} | p{3cm} | p{2cm} | X |} \hline
\rowcolor{gray!33} No. & Date & Version & Approved by \\ \hline

1 & & \arrowversion  &  \\ \hline

\end{tabularx}

\end{document}
